\chapter[26]{Bounds for the Heat Kernel for Evolving Metrics}
\label{ch:chow-iii-heat-kernel-bounds-evolving-metrics}
\dcref{ch26}
\source{chow-et-al-ricci-flow-part-iii:page-0354}

Let \(g(\tau)\), \(\tau\in[0,T]\), satisfy
\[
 \partial_\tau g_{ij}=2\mathcal R_{ij},\qquad
 \mathcal R=g^{ij}\mathcal R_{ij},\qquad
 \partial_\tau d\mu_{g(\tau)}=\mathcal R\,d\mu_{g(\tau)}.
\]
For a smooth potential \(Q\), put
\[
 L=\partial_\tau-\Delta_{g(\tau)}+Q,
 \qquad L^*=-\partial_\tau-\Delta_{g(\tau)}+Q-\mathcal R.
\]
Write \(H(x,\tau;y,v)\) and \(H^*(x,\tau;y,v)\) for the corresponding
heat and adjoint heat kernels.

\section{Heat kernels for evolving metrics}

\begin{lemma}[Duhamel's principle on a closed manifold]
\label{lem:chow-iii-duhamel-evolving-closed}
\source{chow-et-al-ricci-flow-part-iii:page-0355,chow-et-al-ricci-flow-part-iii:page-0356}
\dcref{ch26:26.1}
\group{Evolving Heat Kernels}

If \(A,B\) are \(C^2\) in space and \(C^1\) in time on a closed
evolving manifold, then
\[
 \int_{\tau_1}^{\tau_2}\!\int_{\mathcal M}
 ((LA)B-A(L^*B))\,d\mu_{g(\tau)}d\tau
 =\int_{\mathcal M}AB\,d\mu_{g(\tau_2)}
  -\int_{\mathcal M}AB\,d\mu_{g(\tau_1)}.
\]
\end{lemma}
\begin{proof}
\sketch Expand \(L,L^*\), use the evolution of the volume form, and integrate the
two Laplacian terms by parts.  The integrand becomes
\(\frac d{d\tau}\int_{\mathcal M}AB\,d\mu_{g(\tau)}\).
\end{proof}

\begin{claim}[Duhamel's principle with Dirichlet boundary data]
\label{clm:chow-iii-duhamel-dirichlet-exercise}
\source{chow-et-al-ricci-flow-part-iii:page-0356}
\uses{lem:chow-iii-duhamel-evolving-closed}
\dcref{ch26:26.2}
\group{Evolving Heat Kernels}


Prove the same identity on a compact manifold with boundary when
\(A=B=0\) on \(\partial\mathcal M\times[0,T]\).
\end{claim}

\begin{lemma}[Symmetry of heat and adjoint heat kernels]
\label{lem:chow-iii-heat-adjoint-kernel-symmetry-closed}
\source{chow-et-al-ricci-flow-part-iii:page-0357}
\uses{lem:chow-iii-duhamel-evolving-closed, thm:chow-iii-closed-evolving-heat-kernel}
\dcref{ch26:26.3}
\group{Evolving Heat Kernels}


On a closed manifold, for \(0\le \rho<v\le T\),
\[
 H(y,v;z,\rho)=H^*(z,\rho;y,v).
\]
\end{lemma}
\begin{proof}
\sketch Apply Duhamel's identity to \(H(\cdot,\tau;z,\rho)\) and
\(H^*(\cdot,\tau;y,v)\), then let the two intermediate times tend to
\(\rho\) and \(v\).
\end{proof}

\begin{lemma}[The second variables satisfy the adjoint heat equation]
\label{lem:chow-iii-heat-kernel-second-variables-adjoint}
\source{chow-et-al-ricci-flow-part-iii:page-0357}
\uses{lem:chow-iii-heat-adjoint-kernel-symmetry-closed}
\dcref{ch26:26.4}
\group{Evolving Heat Kernels}


With respect to ((y,v)),
\[
 (\partial_v+\Delta_{y,v}+\mathcal R-Q)H(x,\tau;y,v)=0,
 \qquad \lim_{v\nearrow\tau}H(x,\tau;\cdot,v)=\delta_x.
\]
\end{lemma}

\begin{lemma}[\(L^1\)-norm of the heat kernel on a closed manifold]
\label{lem:chow-iii-closed-heat-kernel-mass-bound}
\source{chow-et-al-ricci-flow-part-iii:page-0357,chow-et-al-ricci-flow-part-iii:page-0358}
\uses{thm:chow-iii-closed-evolving-heat-kernel}
\dcref{ch26:26.5}
\group{Heat Kernel Mass}


If \(C_1=\sup_{\mathcal M\times[0,T]}|\mathcal R-Q|<\infty\), then
\[
 e^{-C_1(\tau-v)}\le
 \int_{\mathcal M}H(x,\tau;y,v)\,d\mu_{g(\tau)}(x)
 \le e^{C_1(\tau-v)}.
\]
\end{lemma}
\begin{proof}
\sketch Differentiate the integral, use \(LH=0\), and apply Gronwall to
\(I'=\int(\mathcal R-Q)H\).
\end{proof}

\begin{corollary}[Mass preservation when \(Q=\mathcal R\)]
\label{cor:chow-iii-closed-heat-kernel-mass-preservation}
\source{chow-et-al-ricci-flow-part-iii:page-0358}
\uses{lem:chow-iii-closed-heat-kernel-mass-bound}
\dcref{ch26:26.6}
\group{Heat Kernel Mass}


If \(Q=\mathcal R\), then
\[
 \int_{\mathcal M}H(x,\tau;y,v)\,d\mu_{g(\tau)}(x)=1.
\]
\end{corollary}

\begin{lemma}[\(L^1\)-norm in the second space variable]
\label{lem:chow-iii-closed-heat-kernel-second-mass-bound}
\source{chow-et-al-ricci-flow-part-iii:page-0358}
\uses{lem:chow-iii-heat-kernel-second-variables-adjoint}
\dcref{ch26:26.7}
\group{Heat Kernel Mass}


If \(C_2=\sup_{\mathcal M\times[0,T]}|Q|\), then
\[
 e^{-C_2(\tau-v)}\le
 \int_{\mathcal M}H(x,\tau;y,v)\,d\mu_{g(v)}(y)
 \le e^{C_2(\tau-v)}.
\]
\end{lemma}

\begin{lemma}[Semigroup property on a closed manifold]
\label{lem:chow-iii-closed-heat-kernel-semigroup}
\source{chow-et-al-ricci-flow-part-iii:page-0359}
\uses{lem:chow-iii-heat-kernel-second-variables-adjoint}
\dcref{ch26:26.8}
\group{Evolving Heat Kernels}


For \(0\le v<\tau<\rho\le T\),
\[
 \int_{\mathcal M}H(z,\rho;x,\tau)H(x,\tau;y,v)
 \,d\mu_{g(\tau)}(x)=H(z,\rho;y,v).
\]
\end{lemma}
\begin{proof}
\sketch The derivative with respect to the intermediate time vanishes by the forward
and adjoint equations.  Let the intermediate time tend to \(v\).
\end{proof}

\begin{lemma}[Symmetry of Dirichlet heat and adjoint kernels]
\label{lem:chow-iii-dirichlet-heat-adjoint-symmetry}
\source{chow-et-al-ricci-flow-part-iii:page-0359,chow-et-al-ricci-flow-part-iii:page-0360}
\uses{clm:chow-iii-duhamel-dirichlet-exercise, thm:chow-iii-evolving-dirichlet-heat-kernel}
\dcref{ch26:26.9}
\group{Dirichlet Heat Kernels}


On a compact manifold with boundary,
\[
 H_D(x,\tau;y,v)=H_D^*(y,v;x,\tau).
\]
In particular \(H_D\) satisfies the adjoint equation and terminal delta
condition in its second pair of variables.
\end{lemma}

\begin{lemma}[\(L^1\)-norm of the Dirichlet heat kernel]
\label{lem:chow-iii-dirichlet-heat-kernel-mass-bounds}
\source{chow-et-al-ricci-flow-part-iii:page-0360,chow-et-al-ricci-flow-part-iii:page-0361,chow-et-al-ricci-flow-part-iii:page-0362}
\uses{lem:chow-iii-dirichlet-heat-adjoint-symmetry}
\dcref{ch26:26.10}
\group{Dirichlet Heat Kernels}


For an interior point \(y\),
\[
 \int_{\mathcal M}H_D(x,\tau;y,v)\,d\mu_{g(\tau)}(x)
 \le e^{C_1(\tau-v)}.
\]
There is also a constant \(C<\infty\), depending only on the uniform metric
evolution, curvature, injectivity, and boundary-distance data specified in
the source, such that
\[
 \int_{\mathcal M}H_D(x,\tau;y,v)\,d\mu_{g(\tau)}(x)
 \ge 1-C(\tau-v).
\]
\end{lemma}
\begin{proof}
\sketch The outward normal derivative of the positive Dirichlet kernel is
nonpositive, giving the upper differential inequality.  For the lower bound,
integrate against a time-dependent cutoff supported away from the boundary;
uniform bounds for its time derivative and Laplacian yield the result.
\end{proof}

\begin{claim}[Alternative lower Dirichlet mass bound]
\label{clm:chow-iii-dirichlet-mass-lower-bound-exercise}
\source{chow-et-al-ricci-flow-part-iii:page-0362}
\uses{lem:chow-iii-dirichlet-heat-kernel-mass-bounds}
\dcref{ch26:26.11}
\group{Dirichlet Heat Kernels}


Derive another lower bound for the Dirichlet heat-kernel mass by the method
used in the derivation of equation (22.41).
\end{claim}

\begin{lemma}[Semigroup property for the Dirichlet heat kernel]
\label{lem:chow-iii-dirichlet-heat-kernel-semigroup}
\source{chow-et-al-ricci-flow-part-iii:page-0362}
\uses{clm:chow-iii-duhamel-dirichlet-exercise}
\dcref{ch26:26.12}
\group{Dirichlet Heat Kernels}


The Dirichlet heat kernel satisfies the semigroup identity of
Lemma~\ref{lem:chow-iii-closed-heat-kernel-semigroup}.
\end{lemma}

\begin{lemma}[Heat-adjoint symmetry on a noncompact manifold]
\label{lem:chow-iii-noncompact-heat-adjoint-symmetry}
\source{chow-et-al-ricci-flow-part-iii:page-0363}
\uses{lem:chow-iii-dirichlet-heat-adjoint-symmetry, thm:chow-iii-noncompact-evolving-heat-kernel}
\dcref{ch26:26.13}
\group{Noncompact Heat Kernels}


For the minimal positive fundamental solution on a complete noncompact
manifold,
\[
 H(x,\tau;y,v)=H^*(y,v;x,\tau).
\]
Thus the second variables satisfy the adjoint equation and terminal delta
condition.
\end{lemma}

\begin{lemma}[\(L^1\)-norm on a complete noncompact manifold]
\label{lem:chow-iii-noncompact-heat-kernel-mass-bound}
\source{chow-et-al-ricci-flow-part-iii:page-0363,chow-et-al-ricci-flow-part-iii:page-0364,chow-et-al-ricci-flow-part-iii:page-0365}
\uses{lem:chow-iii-dirichlet-heat-kernel-mass-bounds, lem:chow-iii-noncompact-heat-adjoint-symmetry}
\dcref{ch26:26.14}
\group{Noncompact Heat Kernels}


Suppose \(g(\tau)\) is complete and noncompact and
\[
 \sup|\sec(g(0))|+\sup|\mathcal R_{ij}|+\sup|\nabla_i\mathcal R_{jk}|<\infty.
\]
If \(C_1=\sup|\mathcal R-Q|\), then
\[
 e^{-C_1(\tau-v)}\le\int_{\mathcal M}H(x,\tau;y,v)
 \,d\mu_{g(\tau)}(x)\le e^{C_1(\tau-v)}.
\]
\end{lemma}
\begin{proof}
\sketch The upper bound follows by exhausting \(\mathcal M\) with Dirichlet kernels.
For the lower bound, integrate against distance cutoffs whose Laplacians tend
uniformly to zero, then let the cutoff radius tend to infinity.
\end{proof}

\begin{corollary}[Mass preservation in the noncompact case]
\label{cor:chow-iii-noncompact-heat-kernel-mass-preservation}
\source{chow-et-al-ricci-flow-part-iii:page-0363}
\uses{lem:chow-iii-noncompact-heat-kernel-mass-bound}
\dcref{ch26:26.15}
\group{Noncompact Heat Kernels}


Under the hypotheses of Lemma~\ref{lem:chow-iii-noncompact-heat-kernel-mass-bound},
if \(Q=\mathcal R\), then
\[
 \int_{\mathcal M}H(x,\tau;y,v)\,d\mu_{g(\tau)}(x)=1.
\]
\end{corollary}

\begin{lemma}[Semigroup property in the noncompact case]
\label{lem:chow-iii-noncompact-heat-kernel-semigroup}
\source{chow-et-al-ricci-flow-part-iii:page-0365,chow-et-al-ricci-flow-part-iii:page-0366}
\uses{lem:chow-iii-dirichlet-heat-kernel-semigroup, thm:chow-iii-noncompact-evolving-heat-kernel}
\dcref{ch26:26.16}
\group{Noncompact Heat Kernels}


For \(0\le v<\rho<\tau\le T\),
\[
 H(x,\tau;y,v)=\int_{\mathcal M}H(x,\tau;z,\rho)
 H(z,\rho;y,v)\,d\mu_{g(\rho)}(z).
\]
\end{lemma}
\begin{proof}
\sketch Use the identity on an exhaustion by Dirichlet domains.  Monotone convergence
gives one inequality; convergence on each compact subset gives the other.
\end{proof}

\section{Upper and lower heat-kernel bounds}

Fix a reference metric \(\widetilde g\) uniformly equivalent to \(g(\tau)\),
as in Chapter~25.

\begin{lemma}[Rough upper bound for \(H\)]
\label{lem:chow-iii-rough-evolving-heat-kernel-upper-bound}
\source{chow-et-al-ricci-flow-part-iii:page-0366,chow-et-al-ricci-flow-part-iii:page-0367,chow-et-al-ricci-flow-part-iii:page-0368}
\uses{lem:chow-iii-noncompact-heat-kernel-mass-bound}
\dcref{ch26:26.17}
\group{Heat Kernel Upper Bounds}


Under the hypotheses of the evolving-metric mean value inequality, there is
\(C<\infty\) such that
\[
 H(x,\tau;y,v)\le\min\left\{
 \frac{C}{\Vol_{\widetilde g}B_{\widetilde g}(x,\sqrt{\tau-v}/2)},
 \frac{C}{\Vol_{\widetilde g}B_{\widetilde g}(y,\sqrt{\tau-v}/2)}
 \right\}.
\]
\end{lemma}
\begin{proof}
\sketch Apply the parabolic mean value inequality at radius \(\sqrt{\tau-v}/2\),
then use the uniform \(L^1\) bound.  Apply the adjoint version for the estimate
in the second variable.
\end{proof}

\begin{remark}[Interpretation of the rough upper bound]
\label{rem:chow-iii-rough-upper-bound-interpretation}
\source{chow-et-al-ricci-flow-part-iii:page-0367}
\uses{lem:chow-iii-rough-evolving-heat-kernel-upper-bound}
\dcref{ch26:26.18}
\group{Heat Kernel Upper Bounds}


The bound is most effective near the diagonal.  For a closed static manifold,
by contrast, the heat kernel tends to \(\Vol(\mathcal M)^{-1}\) as time tends
to infinity.
\end{remark}

\begin{claim}[Rough Dirichlet heat-kernel bounds]
\label{clm:chow-iii-rough-dirichlet-upper-bound-exercise}
\source{chow-et-al-ricci-flow-part-iii:page-0368}
\uses{lem:chow-iii-rough-evolving-heat-kernel-upper-bound}
\dcref{ch26:26.19}
\group{Heat Kernel Upper Bounds}


Prove the analogous bounds for \(H_\Omega\), with a constant independent of
\(\Omega\), whenever the relevant ball of radius \(\sqrt{\tau-v}/2\) lies
inside \(\Omega\).
\end{claim}

\begin{definition}[\((\gamma,A)\)-regular function]
\label{def:chow-iii-gamma-a-regular-function}
\source{chow-et-al-ricci-flow-part-iii:page-0368}
\dcref{ch26:26.20}
\group{Weighted Heat Estimates}

A continuous strictly increasing (f:\(0,T]\to(0,\infty)\) is
\((\gamma,A)\)-regular, where \(\gamma>1\) and \(A\ge1\), if
\[
 \frac{f(\tau_1)}{f(\tau_1/\gamma)}\le
 A\frac{f(\tau_2)}{f(\tau_2/\gamma)}
 \quad(0<\tau_1\le\tau_2<T).
\]
\end{definition}

\begin{lemma}[Exponentially weighted \(L^2\)-estimate]
\label{lem:chow-iii-weighted-l2-heat-estimate}
\source{chow-et-al-ricci-flow-part-iii:page-0369,chow-et-al-ricci-flow-part-iii:page-0370,chow-et-al-ricci-flow-part-iii:page-0371,chow-et-al-ricci-flow-part-iii:page-0372,chow-et-al-ricci-flow-part-iii:page-0373,chow-et-al-ricci-flow-part-iii:page-0374}
\uses{def:chow-iii-gamma-a-regular-function}
\dcref{ch26:26.21}
\group{Weighted Heat Estimates}


Let \(u_\tau=\Delta_{g(\tau)}u-Qu\) on a compact domain \(\Omega\), with
zero Dirichlet data and initial support in a compact \(\mathcal K\Subset\Omega\).
If
\[
 \int_\Omega u^2(x,\tau)\,d\mu_{g(\tau)}(x)\le f(\tau)^{-1}
\]
for a \((\gamma,A)\)-regular \(f\), then constants \(C,D>0\), with the
dependence specified in the source, satisfy
\[
 \int_\Omega u^2(x,\tau)
 \exp\!\left(\frac{d_{g(0)}(x,\mathcal K)^2}{D\tau}\right)
 d\mu_{g(\tau)}(x)
 \le \frac{4Ae^{C\tau}}{f(\tau/\gamma)}.
\]
\end{lemma}
\begin{proof}
\sketch After replacing (u) by \(e^{-C_0\tau}u\), multiply its equation by a
Davies weight \(e^\xi u\), where
\(\xi=-\rho_R^2/(2\widetilde C_0(\sigma-\tau))\).  The inequality
\(\xi_\tau+|\nabla\xi|^2/2\le0\) makes the weighted \(L^2\) norm
nonincreasing.  Iterating the resulting tail estimate at times
\(\tau/\gamma^i\) and radii tending to \(R/2\), using regularity of \(f\),
gives Gaussian decay of the tails.  Summing those tails over annuli proves
the displayed estimate.
\end{proof}

\begin{remark}[Point-supported initial data]
\label{rem:chow-iii-weighted-l2-point-source}
\source{chow-et-al-ricci-flow-part-iii:page-0369}
\uses{lem:chow-iii-weighted-l2-heat-estimate}
\dcref{ch26:26.22}
\group{Weighted Heat Estimates}


The important special case is \(\mathcal K=\{y\}\), with (u) a Dirichlet
heat kernel based at \(y\).
\end{remark}

\begin{lemma}[Average exponential quadratic decay]
\label{lem:chow-iii-heat-kernel-average-gaussian-decay}
\source{chow-et-al-ricci-flow-part-iii:page-0375,chow-et-al-ricci-flow-part-iii:page-0376}
\uses{lem:chow-iii-rough-evolving-heat-kernel-upper-bound, lem:chow-iii-weighted-l2-heat-estimate}
\dcref{ch26:26.23}
\group{Heat Kernel Upper Bounds}


If \(\Ric_{g(0)}\ge-K\), then constants \(C,D<\infty\) satisfy
\[
 \int_{\mathcal M}H^2(x,\tau;y,v)
 e^{d_{g(0)}(x,y)^2/(D(\tau-v))}\,d\mu_{g(\tau)}(x)
 \le\frac{C}{\Vol_{g(0)}B_{g(0)}(y,\sqrt{\tau-v}/4)},
\]
and the analogous estimate obtained by integrating in \(y\) with measure
\(d\mu_{g(v)}\) and the volume centered at \(x\).
\end{lemma}
\begin{proof}
\sketch Apply Lemma~\ref{lem:chow-iii-weighted-l2-heat-estimate} to Dirichlet kernels
on an exhaustion.  The rough pointwise bound and the mass bound provide the
unweighted \(L^2\) hypothesis.  Bishop--Gromov comparison makes the ball-volume
function ((4,A))-regular.  Pass to the exhaustion limit.
\end{proof}

\begin{claim}[Point source in the weighted estimate]
\label{clm:chow-iii-weighted-l2-delta-source-exercise}
\source{chow-et-al-ricci-flow-part-iii:page-0376}
\uses{lem:chow-iii-weighted-l2-heat-estimate}
\dcref{ch26:26.24}
\group{Weighted Heat Estimates}


Justify the application of Lemma~\ref{lem:chow-iii-weighted-l2-heat-estimate}
when the initial datum is a delta distribution.
\end{claim}

\begin{theorem}[Upper bound for the heat kernel for an evolving metric]
\label{thm:chow-iii-evolving-heat-kernel-gaussian-upper-bound}
\source{chow-et-al-ricci-flow-part-iii:page-0376,chow-et-al-ricci-flow-part-iii:page-0377}
\uses{lem:chow-iii-heat-kernel-average-gaussian-decay, lem:chow-iii-noncompact-heat-kernel-semigroup}
\dcref{ch26:26.25}
\group{Heat Kernel Upper Bounds}


There are constants \(C_3,C_4<\infty\), with the dependence specified in the
source, such that
\[
 H(x,\tau;y,v)\le
 \frac{C_3e^{-d_{\widetilde g}(x,y)^2/(C_4(\tau-v))}}
 {\Vol_{\widetilde g}^{1/2}B_{\widetilde g}(x,\sqrt{(\tau-v)/2})
  \Vol_{\widetilde g}^{1/2}B_{\widetilde g}(y,\sqrt{(\tau-v)/2})}.
\]
\end{theorem}
\begin{proof}
\sketch Split the semigroup identity at \(\rho=(\tau+v)/2\), distribute half the
Gaussian distance between its two factors, and apply Cauchy--Schwarz and the
two weighted \(L^2\) estimates.
\end{proof}

\begin{corollary}[One-sided volume Gaussian upper bound]
\label{cor:chow-iii-evolving-heat-kernel-one-volume-upper-bound}
\source{chow-et-al-ricci-flow-part-iii:page-0377,chow-et-al-ricci-flow-part-iii:page-0378,chow-et-al-ricci-flow-part-iii:page-0379}
\uses{lem:chow-iii-rough-evolving-heat-kernel-upper-bound, thm:chow-iii-evolving-heat-kernel-gaussian-upper-bound}
\dcref{ch26:26.26}
\group{Heat Kernel Upper Bounds}


For constants \(C,D<\infty\),
\[
 H(x,\tau;y,v)\le Ce^{-d_{\widetilde g}(x,y)^2/(D(\tau-v))}
 \min\left\{
 \frac1{\Vol_{\widetilde g}B_{\widetilde g}(x,\sqrt{(\tau-v)/2})},
 \frac1{\Vol_{\widetilde g}B_{\widetilde g}(y,\sqrt{(\tau-v)/2})}
 \right\}.
\]
\end{corollary}
\begin{proof}
\sketch Near the diagonal use the rough bound.  Off the diagonal use the two-volume
bound and Bishop--Gromov comparison to absorb the volume ratio into a weaker
Gaussian.
\end{proof}

\begin{claim}[Space-form volume absorption]
\label{clm:chow-iii-space-form-volume-absorption-exercise}
\source{chow-et-al-ricci-flow-part-iii:page-0379}
\uses{cor:chow-iii-evolving-heat-kernel-one-volume-upper-bound}
\dcref{ch26:26.27}
\group{Heat Kernel Upper Bounds}


Verify the space-form volume estimate (26.68) used to absorb a relative
volume ratio into the Gaussian factor.
\end{claim}

\begin{lemma}[\(L^1\)-mass concentrates at the pole]
\label{lem:chow-iii-heat-kernel-mass-concentration}
\source{chow-et-al-ricci-flow-part-iii:page-0379,chow-et-al-ricci-flow-part-iii:page-0380}
\uses{cor:chow-iii-evolving-heat-kernel-one-volume-upper-bound}
\dcref{ch26:26.28}
\group{Heat Kernel Lower Bounds}


Uniformly for \(0\le v<\tau\le T\),
\[
 \int_{\mathcal M\setminus B_{\widetilde g}(y,A\sqrt{\tau-v})}
 H(x,\tau;y,v)\,d\mu_{g(\tau)}(x)\longrightarrow0
\]
as \(A\to\infty\), and likewise after interchanging the two space variables
and integrating with \(d\mu_{g(v)}\).
\end{lemma}
\begin{proof}
\sketch Integrate the one-volume Gaussian upper bound in polar shells.  Relative
volume comparison gives at most exponential linear growth, which is dominated
uniformly by the Gaussian tail.
\end{proof}

\begin{corollary}[Positive mass in a parabolic-scale ball]
\label{cor:chow-iii-heat-kernel-local-mass-lower-bound}
\source{chow-et-al-ricci-flow-part-iii:page-0380}
\uses{lem:chow-iii-heat-kernel-mass-concentration, lem:chow-iii-noncompact-heat-kernel-mass-bound}
\dcref{ch26:26.29}
\group{Heat Kernel Lower Bounds}


There are \(A<\infty\) and \(c>0\) such that
\[
 \int_{B_{\widetilde g}(y,A\sqrt{\tau-v})}H(x,\tau;y,v)
 \,d\mu_{g(\tau)}(x)\ge c,
\]
and analogously in the second space variable.
\end{corollary}

\begin{lemma}[Lower bound along the diagonal]
\label{lem:chow-iii-evolving-heat-kernel-diagonal-lower-bound}
\source{chow-et-al-ricci-flow-part-iii:page-0380,chow-et-al-ricci-flow-part-iii:page-0381}
\uses{cor:chow-iii-heat-kernel-local-mass-lower-bound}
\dcref{ch26:26.30}
\group{Heat Kernel Lower Bounds}


There exists \(c>0\) such that
\[
 H(y,\tau;y,v)\ge
 \frac{c}{\Vol_{\widetilde g}B_{\widetilde g}(y,\sqrt{\tau-v})}.
\]
\end{lemma}
\begin{proof}
\sketch The Li--Yau Harnack inequality compares the diagonal value at time \(\tau\)
with values on a parabolic-scale ball at the midpoint time.  Integrate that
comparison and apply Corollary~\ref{cor:chow-iii-heat-kernel-local-mass-lower-bound}.
\end{proof}

\begin{theorem}[Lower bound for the heat kernel]
\label{thm:chow-iii-evolving-heat-kernel-gaussian-lower-bound}
\source{chow-et-al-ricci-flow-part-iii:page-0381,chow-et-al-ricci-flow-part-iii:page-0382}
\uses{lem:chow-iii-evolving-heat-kernel-diagonal-lower-bound}
\dcref{ch26:26.31}
\group{Heat Kernel Lower Bounds}


There are \(c_1,c_2>0\), with the curvature dependence specified in the
source, such that
\[
 H(x,\tau;y,v)\ge c_1
 \min\left\{
 \frac1{\Vol_{\widetilde g}B_{\widetilde g}(x,\sqrt{\tau-v})},
 \frac1{\Vol_{\widetilde g}B_{\widetilde g}(y,\sqrt{\tau-v})}
 \right\}
 e^{-d_{\widetilde g}(x,y)^2/(c_2(\tau-v))}.
\]
Equivalently, after changing constants, one may use the geometric mean of
the two reciprocal volumes.
\end{theorem}
\begin{proof}
\sketch Apply the Li--Yau Harnack inequality between \((x,\tau)\) and the diagonal
point (\(y,(\tau+v)/2\)), then insert the diagonal lower bound.  Repeat with
(x,y) interchanged and combine the two estimates.
\end{proof}

\begin{theorem}[Davies' double integral upper bound]
\label{thm:chow-iii-davies-double-integral-upper-bound}
\source{chow-et-al-ricci-flow-part-iii:page-0382,chow-et-al-ricci-flow-part-iii:page-0383,chow-et-al-ricci-flow-part-iii:page-0384}
\uses{lem:chow-iii-noncompact-heat-kernel-semigroup, lem:chow-iii-weighted-l2-heat-estimate}
\dcref{ch26:26.32}
\group{Heat Kernel Upper Bounds}


For a complete static Riemannian manifold, its heat kernel satisfies
\[
 \int_{U_1}\!\int_{U_2}H(x,y,t)\,d\mu(x)d\mu(y)
 \le \Vol(U_1)^{1/2}\Vol(U_2)^{1/2}
 e^{-d(U_1,U_2)^2/(4t)}
\]
for bounded \(U_1,U_2\subset\mathcal M\).
\end{theorem}
\begin{proof}
\sketch Split the kernel at (t/2).  For
\(u_i(z,t)=\int_{U_i}H(z,y,t)\,d\mu(y)\), use Cauchy--Schwarz with the
weights \(e^{d(z,U_i)^2/(2t)}\).  Their weighted \(L^2\) norms decrease by
the Davies monotonicity calculation and converge at time zero to
\(\Vol(U_i)\).  The triangle inequality supplies the displayed Gaussian.
\end{proof}

\begin{remark}[Pointwise bounds from Davies' estimate]
\label{rem:chow-iii-davies-estimate-pointwise-use}
\source{chow-et-al-ricci-flow-part-iii:page-0384}
\uses{thm:chow-iii-davies-double-integral-upper-bound}
\dcref{ch26:26.33}
\group{Heat Kernel Upper Bounds}


Davies' estimate gives an alternative proof of pointwise upper bounds on
complete manifolds satisfying a parabolic mean value inequality.
\end{remark}

\section{Heat balls and space-time mean value properties}

\begin{theorem}[Mean value property for harmonic functions on \(\mathbb E^n\)]
\label{thm:chow-iii-euclidean-harmonic-mean-value-property}
\source{chow-et-al-ricci-flow-part-iii:page-0384,chow-et-al-ricci-flow-part-iii:page-0385}
\dcref{ch26:26.35}
\group{Mean Value Properties}

If \(\Delta u=0\) in \(\Omega\subset\mathbb R^n\), then
\[
 u(O)=\frac1{\omega_nr^n}\int_{B(O,r)}u\,d\mu_E
\]
whenever \(B(O,r)\subset\Omega\).
\end{theorem}
\begin{proof}
\sketch Differentiate the average in (r).  Two integrations by parts using
\(\Delta|x|^2=2n\) show that its derivative vanishes; its limit at zero is
(u(O)).
\end{proof}

\begin{claim}[Mean value inequality when \(\Ric\ge0\)]
\label{clm:chow-iii-nonpositive-subharmonic-mean-value-exercise}
\source{chow-et-al-ricci-flow-part-iii:page-0386}
\dcref{ch26:26.36}
\group{Mean Value Properties}

If \(\Ric\ge0\), \(f\le0\), and \(\Delta f\ge0\), prove that
\[
 f(O)\le\frac1{\omega_nr^n}\int_{B(O,r)}f\,d\mu
\]
for \(0<r<\operatorname{inj}(O)\).
\end{claim}

\begin{remark}[Euclidean volume upper bound]
\label{rem:chow-iii-nonnegative-ricci-euclidean-volume-bound}
\source{chow-et-al-ricci-flow-part-iii:page-0386}
\uses{clm:chow-iii-nonpositive-subharmonic-mean-value-exercise}
\dcref{ch26:26.37}
\group{Mean Value Properties}


Taking \(f=-1\) gives \(\Vol B(O,r)\le\omega_nr^n\), the basic
Bishop--Gromov comparison.
\end{remark}

For the Euclidean heat kernel \(H_E\), define the heat ball and its logarithmic
weight by
\[
 E_r(x,t)=\{(y,s):H_E(x,t;y,s)>r^{-n}\},
 \qquad \psi_r=\log(r^nH_E).
\]

\begin{claim}[Normalization of the Euclidean heat-ball weight]
\label{clm:chow-iii-euclidean-heat-ball-normalization-exercise}
\source{chow-et-al-ricci-flow-part-iii:page-0387}
\dcref{ch26:26.38}
\group{Heat Balls}

Prove
\[
 \frac1{4r^n}\iint_{E_r(x,t)}\frac{|x-y|^2}{(t-s)^2}
 \,d\mu_E(y)ds=1.
\]
\end{claim}

\begin{theorem}[Space-time mean value property on Euclidean space]
\label{thm:chow-iii-euclidean-heat-mean-value-property}
\source{chow-et-al-ricci-flow-part-iii:page-0387,chow-et-al-ricci-flow-part-iii:page-0388}
\uses{clm:chow-iii-euclidean-heat-ball-normalization-exercise, lem:chow-iii-euclidean-heat-ball-average-stationary}
\dcref{ch26:26.39}
\group{Heat Balls}


If \(u_t=\Delta u\) on \(\mathbb R^n\times[0,t]\) and
\(t\ge r^2/(4\pi)>0\), then
\[
 u(x,t)=\frac1{4r^n}\iint_{E_r(x,t)}u(y,s)
 \frac{|x-y|^2}{(t-s)^2}\,d\mu_E(y)ds.
\]
\end{theorem}

\begin{lemma}[The Euclidean heat-ball average is stationary]
\label{lem:chow-iii-euclidean-heat-ball-average-stationary}
\source{chow-et-al-ricci-flow-part-iii:page-0388}
\dcref{ch26:26.40}
\group{Heat Balls}

For a heat solution, the weighted average on the right side of
Theorem~\ref{thm:chow-iii-euclidean-heat-mean-value-property} has zero
derivative with respect to (r).
\end{lemma}

On a complete static manifold, define
\[
 E_r(x_0,t_0)=\{(y,t):H(x_0,y,t_0-t)>r^{-n}\},
 \qquad \psi_r=\log H+n\log r.
\]

\begin{theorem}[Space-time mean value property on a manifold]
\label{thm:chow-iii-riemannian-heat-mean-value-property}
\source{chow-et-al-ricci-flow-part-iii:page-0390,chow-et-al-ricci-flow-part-iii:page-0391,chow-et-al-ricci-flow-part-iii:page-0392,chow-et-al-ricci-flow-part-iii:page-0393,chow-et-al-ricci-flow-part-iii:page-0394}
\uses{def:chow-iii-gamma-a-regular-function}
\dcref{ch26:26.41}
\group{Heat Balls}


Suppose \(H(y,y,\tau)\le C/f(\tau)\) for a \((\gamma,A)\)-regular
\(f:(0,\infty)\to(0,\infty)\).  If \(u_t=\Delta u\), then for sufficiently
small \(r>0\), \(E_r(x_0,t_0)\) is relatively compact and
\[
 u(x_0,t_0)=\frac1{r^n}\iint_{E_r(x_0,t_0)}u(y,t)
 |\nabla\log H(x_0,y,t_0-t)|^2\,d\mu(y)dt.
\]
\end{theorem}
\begin{proof}
\sketch The diagonal hypothesis yields a Gaussian off-diagonal bound, hence compact
heat-ball slices and finite integrals.  The coarea formula and
\(\psi_{r,t}+\Delta\psi_r=-|\nabla\psi_r|^2\) give
\[
 \frac d{dr}I(r)=\frac n{r^{n+1}}\iint_{E_r}\psi_r(\Delta u-u_t)\,d\mu dt.
\]
Thus \(I\) is constant for heat solutions, and heat-kernel asymptotics give
\(\lim_{r\downarrow0}I(r)=u(x_0,t_0)\).
\end{proof}

\begin{theorem}[Mean value property using heat spheres]
\label{thm:chow-iii-heat-sphere-mean-value-property}
\source{chow-et-al-ricci-flow-part-iii:page-0395}
\uses{thm:chow-iii-riemannian-heat-mean-value-property}
\dcref{ch26:26.42}
\group{Heat Balls}


Under the hypotheses of Theorem~\ref{thm:chow-iii-riemannian-heat-mean-value-property},
\[
 u(x_0,0)=\int_{\partial E_r}u(y,t)
 \frac{|\nabla H|^2}{\sqrt{|\nabla H|^2+|\partial_tH|^2}}
 \,d\widetilde\sigma(y,t).
\]
\end{theorem}
\begin{proof}
\sketch Apply the space-time divergence theorem to
\(uH\partial_t+u\nabla H-H\nabla u\) on truncated heat balls and let the
top time tend to zero.
\end{proof}

\begin{claim}[Recover the heat-ball formula from heat spheres]
\label{clm:chow-iii-heat-sphere-coarea-exercise}
\source{chow-et-al-ricci-flow-part-iii:page-0395}
\uses{thm:chow-iii-heat-sphere-mean-value-property, thm:chow-iii-riemannian-heat-mean-value-property}
\dcref{ch26:26.43}
\group{Heat Balls}


Derive the heat-ball mean value formula by integrating the heat-sphere formula
and applying the coarea formula.
\end{claim}

For a Ricci flow and a basepoint \((x_0,t_0)\), let \(\ell(y,\tau)\) be the
reduced distance, \(v=(4\pi\tau)^{-n/2}e^{-\ell}\), and set
\[
 E_r=\{(y,t):v(y,t_0-t)\ge r^{-n}\},\quad
 \psi_r=\log v+n\log r,
\]
\[
 P(r)=\iint_{E_r}(|\nabla\psi_r|^2+\psi_rR)u\,d\mu_{g(t)}dt.
\]

\begin{theorem}[Local monotonicity formula under Ricci flow]
\label{thm:chow-iii-ricci-flow-local-mean-value-monotonicity}
\source{chow-et-al-ricci-flow-part-iii:page-0396,chow-et-al-ricci-flow-part-iii:page-0397}
\dcref{ch26:26.44}
\group{Ricci Flow Heat Balls}

For \(I(r)=P(r)/r^n\) and \(0<r_1<r_2\),
\[
 I(r_2)-I(r_1)=-\int_{r_1}^{r_2}\frac n{r^{n+1}}
 \iint_{E_r}\left[
 \frac uv(\partial_t+\Delta-R)v+\psi_r(\partial_t-\Delta)u
 \right]d\mu_{g(t)}dt\,dr.
\]
\end{theorem}

\begin{corollary}[Mean value inequality for heat supersolutions]
\label{cor:chow-iii-ricci-flow-supersolution-mean-value}
\source{chow-et-al-ricci-flow-part-iii:page-0397}
\uses{thm:chow-iii-ricci-flow-local-mean-value-monotonicity}
\dcref{ch26:26.45}
\group{Ricci Flow Heat Balls}


If \(u\ge0\) and \((\partial_t-\Delta)u\ge0\), then (I(r)) is
nonincreasing and
\[
 u(x_0,t_0)\ge\frac1{r^n}\iint_{E_r}
 (|\nabla\psi_r|^2+\psi_rR)u\,d\mu_{g(t)}dt.
\]
\end{corollary}

\begin{corollary}[Strong maximum principle under Ricci flow]
\label{cor:chow-iii-ricci-flow-strong-maximum-principle}
\source{chow-et-al-ricci-flow-part-iii:page-0397}
\uses{cor:chow-iii-ricci-flow-supersolution-mean-value}
\dcref{ch26:26.46}
\group{Ricci Flow Heat Balls}


If \(R\ge0\) and a heat supersolution (w) attains its space-time minimum at
\((x_0,t_0)\), then \(w(x,t)=w(x_0,t_0)\) for all \(t\le t_0\).
\end{corollary}

\begin{theorem}[Weak Harnack inequality]
\label{thm:chow-iii-weak-parabolic-harnack-inequality}
\source{chow-et-al-ricci-flow-part-iii:page-0397,chow-et-al-ricci-flow-part-iii:page-0398}
\dcref{ch26:26.47}
\group{Weak Heat Solutions}

Assume the local \(L^2\) Poincare inequality (26.124) and the local volume
doubling condition (26.125).  Then there are \(p_0>0\) and \(A<\infty\)
such that every positive weak heat supersolution satisfies
\[
 \left(\frac1{\Vol(Q'_-)}\int_{Q'_-}u^{p_0}\,d\mu dt\right)^{1/p_0}
 \le A\inf_{Q_+}u
\]
on the parabolic subregions specified in the source.
\end{theorem}

\begin{theorem}[Harnack estimate]
\label{thm:chow-iii-weak-solution-parabolic-harnack-estimate}
\source{chow-et-al-ricci-flow-part-iii:page-0398}
\uses{thm:chow-iii-weak-parabolic-harnack-inequality}
\dcref{ch26:26.48}
\group{Weak Heat Solutions}


Under the hypotheses of Theorem~\ref{thm:chow-iii-weak-parabolic-harnack-inequality},
every positive weak heat solution satisfies
\[
 \sup_{Q'_-}u\le A\inf_{Q_+}u.
\]
\end{theorem}

\section{Distance-like functions on complete noncompact manifolds}

\begin{proposition}[Distance-like functions with bounds on derivatives]
\label{prop:chow-iii-distance-like-function-bounded-derivatives}
\source{chow-et-al-ricci-flow-part-iii:page-0399,chow-et-al-ricci-flow-part-iii:page-0400,chow-et-al-ricci-flow-part-iii:page-0401,chow-et-al-ricci-flow-part-iii:page-0402,chow-et-al-ricci-flow-part-iii:page-0403,chow-et-al-ricci-flow-part-iii:page-0404,chow-et-al-ricci-flow-part-iii:page-0405,chow-et-al-ricci-flow-part-iii:page-0406}
\uses{cor:chow-iii-noncompact-heat-kernel-mass-preservation, thm:chow-iii-weak-parabolic-harnack-inequality}
\dcref{ch26:26.49}
\group{Distance-like Functions}


For \(n\ge2\) and \(K>0\), there is \(C_{n,K}>1\) such that every complete
noncompact \((\mathcal M^n,g)\) with \( |\sec(g)|\le K\), and every
\(O\in\mathcal M\), admit a smooth \(f:\mathcal M\to\mathbb R\) satisfying
\[
 d(x,O)+1\le f(x)\le d(x,O)+C_{n,K},
 \qquad |\nabla f|+|\Hess f|\le C_{n,K}.
\]
\end{proposition}
\begin{proof}
\sketch Start from a smooth (u) with \(|u-d(\cdot,O)|\le1\) and
\(|\nabla u|\le2\), and evolve it by the static heat kernel.  Gaussian
heat-kernel bounds and volume comparison show that the evolved function stays
within a uniform constant of (u).  The Karp--Li maximum principle applied
to \(e^{-2(n-1)Kt}|\nabla f|^2\), justified by a Gaussian weighted energy
bound, controls the gradient.  Pull back by the exponential map and use
uniform harmonic coordinates and interior parabolic Schauder estimates to
control the Hessian at time one.  Adding a constant yields the stated
distance inequalities.
\end{proof}

\begin{remark}[Bounds for all higher derivatives]
\label{rem:chow-iii-distance-like-function-higher-derivatives}
\source{chow-et-al-ricci-flow-part-iii:page-0399,chow-et-al-ricci-flow-part-iii:page-0406}
\uses{prop:chow-iii-distance-like-function-bounded-derivatives}
\dcref{ch26:26.50}
\group{Distance-like Functions}


If \(|\nabla^k\Rm|\le K\) for \(0\le k\le m\), the function may be chosen so
that
\[
 |\nabla^kf|\le C_{n,K,m}\qquad(1\le k\le m+2).
\]
\end{remark}

\begin{remark}[Uniform convergence of the heat evolution]
\label{rem:chow-iii-distance-like-heat-evolution-uniform-convergence}
\source{chow-et-al-ricci-flow-part-iii:page-0399}
\uses{prop:chow-iii-distance-like-function-bounded-derivatives}
\dcref{ch26:26.51}
\group{Distance-like Functions}


The heat evolution of the initial distance-like function converges uniformly
to it as \(t\downarrow0\); this follows from the local estimate and the
Gaussian tail bounds used in the proof.
\end{remark}

\begin{remark}[Small-time Gaussian tail]
\label{rem:chow-iii-distance-like-small-time-tail}
\source{chow-et-al-ricci-flow-part-iii:page-0402}
\uses{rem:chow-iii-distance-like-heat-evolution-uniform-convergence}
\dcref{ch26:26.52}
\group{Distance-like Functions}


When \(0<t\le\varepsilon\), the Gaussian tail constant in the proof tends to
zero as \(\varepsilon\downarrow0\).
\end{remark}

\begin{remark}[Geodesic and harmonic coordinate regularity]
\label{rem:chow-iii-geodesic-harmonic-coordinate-regularity}
\source{chow-et-al-ricci-flow-part-iii:page-0406}
\uses{rem:chow-iii-distance-like-function-higher-derivatives}
\dcref{ch26:26.53}
\group{Distance-like Functions}


Curvature-derivative bounds control metric derivatives in geodesic
coordinates, while harmonic coordinates realize the maximal regularity used
by the Schauder argument.
\end{remark}

\begin{claim}[Ricci-flow Davies estimate problem]
\label{clm:chow-iii-ricci-flow-davies-estimate-problem}
\dcref{ch26:26.34}
\group{Heat Kernel Upper Bounds}
\source{chow-et-al-ricci-flow-part-iii:page-0384}
\uses{thm:chow-iii-davies-double-integral-upper-bound}
Find a Ricci-flow analogue of Davies' double integral estimate.
\end{claim}
